% TF.tex - Talent Factory page layout, fonts and colors for documentation
% ConTeXt MkIV / LuaMetaTeX

% =============================================================================
% LANGUAGE (German)
% =============================================================================

\mainlanguage[de]
\setuplabeltext[de][content=Inhaltsverzeichnis]
\setuplabeltext[de][figure=Abbildung~]
\setuplabeltext[de][table=Tabelle~]
\setuplabeltext[de][chapter=Kapitel~]  % ~ = non-breaking space after 'Kapitel'

% =============================================================================
% PAGE LAYOUT
% =============================================================================

\setuppapersize[A4]

\setuplayout[
    width=middle,
    height=middle,
    topspace=1.8cm,
    backspace=2.5cm,
    header=1cm,
    footer=1cm,
    headerdistance=0.5cm,
    footerdistance=0.5cm,
]

% =============================================================================
% FONTS
% =============================================================================

% Sans-serif for headings, serif for body text, mono for code
% Heros lacks accented Greek and IPA glyphs; take them from Libertinus Sans.
% A fallback must be defined before the family it extends.
\definefallbackfamily[mainface][sans][Libertinus Sans]
  [range={greekandcoptic,greekextended,ipaextensions,spacingmodifierletters},force=yes]
\definefontfamily[mainface][sans][TeX Gyre Heros]
\definefontfamily[mainface][serif][TeX Gyre Pagella]
\definefontfamily[mainface][mono][TeX Gyre Cursor][scale=1.0]

\setupbodyfont[mainface, 11pt]

% Plain \starttyping blocks match the size of the code environments below
\setuptyping[style=\tt\small]

% =============================================================================
% COLORS
% =============================================================================

% Talent Factory Corporate Identity (source: website/tailwind.config.js)
\definecolor[tfdark][h=0D0D1A]            % tf-dark
\definecolor[tfaccent][h=7C3AED]          % tf-accent (violet)
\definecolor[tfglow][h=A855F7]            % tf-accent-glow
\definecolor[tfparticles][h=818CF8]       % tf-particles
\definecolor[tftext][h=0F172A]            % tf-text
\definecolor[tfmuted][h=64748B]           % tf-muted
\definecolor[tfborder][h=E2E8F0]          % tf-border
\definecolor[tfsurface][h=F8FAFC]         % tf-surface
\definecolor[tfsurfacealt][h=EEF2FF]      % tf-surface-alt

% Semantic roles (documents may override them)
\definecolor[primarycolor][tfaccent]
\definecolor[secondarycolor][h=059669]    % green (tip)
\definecolor[accentcolor][h=D97706]       % orange (accent)
\definecolor[warningcolor][h=DC2626]      % red (warning)
\definecolor[mutedcolor][tfmuted]
\definecolor[codebackground][tfsurface]
\definecolor[codeborder][tfborder]
\definecolor[infobackground][tfsurfacealt]

% =============================================================================
% HEADINGS
% =============================================================================

% Chapter opening: large number in muted color, title below, rule in the primary color
\define[2]\TFchaptercommand{%
    \vbox{%
        \hsize\textwidth
        {\tfd\bf\color[tfmuted]{#1}}\par
        \blank[small]
        {\bfd\bf\color[primarycolor]{#2}}\par
        \blank[small]
        \blackrule[width=\textwidth,height=1.5pt,color=primarycolor]%
    }%
}

\setuphead[chapter][
    style=\bfd\bf,
    color=primarycolor,
    command=\TFchaptercommand,
    before={\blank[2*big]},
    after={\blank[2*big]},
]

\setuphead[section][
    style=\bfb\bf,
    color=primarycolor,
    before={\blank[2*big]},
    after={\blank[big]},
]

\setuphead[subsection][
    style=\bfa\bf,
    color=primarycolor,
    before={\blank[big]},
    after={\blank[medium]},
]

\setuphead[subsubsection][
    style=\bf,
    before={\blank[medium]},
    after={\blank[small]},
]

% =============================================================================
% GIT INFORMATION
% =============================================================================

% Determine git branch and commit id at build time
\startluacode
    local function get_git_info()
        local branch = "unknown"
        local commit = "000000"

        -- Branch
        local handle = io.popen("git rev-parse --abbrev-ref HEAD 2>/dev/null")
        if handle then
            branch = handle:read("*a"):gsub("%s+", "")
            handle:close()
        end

        -- Short commit id
        handle = io.popen("git rev-parse --short HEAD 2>/dev/null")
        if handle then
            commit = handle:read("*a"):gsub("%s+", "")
            handle:close()
        end

        -- Date of the commit in German long form, e.g. "18. Oktober 2026"
        local date = "unbekannt"
        handle = io.popen("git log -1 --format=%cd --date=format:%Y-%m-%d 2>/dev/null")
        if handle then
            local y, m, d = handle:read("*a"):match("(%d+)-(%d+)-(%d+)")
            handle:close()
            if y then
                local months = {
                    "Januar", "Februar", "März", "April", "Mai", "Juni", "Juli",
                    "August", "September", "Oktober", "November", "Dezember",
                }
                date = tonumber(d) .. ". " .. months[tonumber(m)] .. " " .. y
            end
        end

        return branch, commit, date
    end

    -- Short working-copy commit of the Stemma repository in the parent directory (ADR-010)
    local function get_stemma_commit()
        local commit = "unknown"
        local file = io.open("../.stemma/working-copy", "r")
        if file then
            commit = (file:read("*a") or ""):match("b3:(%x+)") or commit
            file:close()
        end
        return commit:sub(1, 12)
    end

    local branch, commit, date = get_git_info()
    local stemma = get_stemma_commit()
    document.gitbranch = branch
    document.gitcommit = commit
    document.gitdate = date
    document.stemmacommit = stemma
    document.gitinfo = "[git] • Branch: " .. branch .. "@" .. commit
    document.stemmainfo = "[stemma] • Working-Copy-Commit: " .. stemma
\stopluacode

% Make the git information available as TeX macros
\def\gitinfo{\ctxlua{context(document.gitinfo)}}
\def\stemmainfo{\ctxlua{context(document.stemmainfo)}}
\def\gitbranch{\ctxlua{context(document.gitbranch)}}
\def\gitcommit{\ctxlua{context(document.gitcommit)}}
\def\gitdate{\ctxlua{context(document.gitdate)}}

% =============================================================================
% HEADERS AND FOOTERS
% =============================================================================

% Enable markings for chapter and section
\setuphead[chapter][marking=chapter]
\setuphead[section][marking=section]

% Page number only in the footer, not in the header
\setuppagenumbering[location=]

% Header: left = short title (variable document:shorttitle, set by the document),
% right = chapter (uppercase)
\setvariables[document][shorttitle={}]
% The rule below the header is part of the header content
\setupheadertexts
    [\vbox{\hbox to \hsize{\getvariable{document}{shorttitle}\hfil\WORD{\getmarking[chapter]}}%
        \kern 0.4em
        \blackrule[width=\hsize,height=0.6pt,color=tfmuted]}]

% Footer: git and Stemma references centered, one above the other; page number right
\setupfootertexts
    [][\vbox{%
        \hbox to \hsize{\hfil\gitinfo\hfil\llap{\pagenumber}}%
        \hbox to \hsize{\hfil\stemmainfo\hfil}}]

\setupheader[style=\small, color=black]
\setupfooter[style=\small, color=black]

% =============================================================================
% PARAGRAPHS AND SPACING
% =============================================================================

\setupwhitespace[medium]
\setupindenting[no]

% =============================================================================
% LISTS
% =============================================================================

% Unicode bullets do not work here: keep the default ConTeXt symbols

\setupitemize[1][
    margin=1em,
    distance=0.5em,
    before={\blank[small]},
    after={\blank[small]},
    inbetween={\blank[none]},
]

\setupitemize[2][
    margin=1.5em,
    distance=0.5em,
    inbetween={\blank[none]},
]

\setupitemize[3][
    margin=2em,
    distance=0.5em,
    inbetween={\blank[none]},
]

% =============================================================================
% LINKS
% =============================================================================

\setupinteraction[
    state=start,
    color=primarycolor,
    contrastcolor=primarycolor,
    style=normal,
]

% =============================================================================
% TABLE OF CONTENTS
% =============================================================================

\setupcombinedlist[content][
    list={chapter,section},
    alternative=c,
]

% Second level starts flush with the title of the chapter line
\setuplist[chapter][
    style=\bf,
    width=2em,
    before={\blank[medium]},
]

\setuplist[section][
    margin=2em,
    width=2.5em,
]

% Tighter line spacing: use \placecompactcontent instead of \placecontent
\define\placecompactcontent{%
    \begingroup
    \setupwhitespace[none]
    \setupinterlinespace[line=1.5em]
    \placecontent
    \endgroup
}

% =============================================================================
% TABLES AND CAPTIONS
% =============================================================================

% Caption above the table: "Tabelle 1: Title"
\setupcaption[table][
    location=top,
    way=bytext,
    prefix=no,
    headstyle=\bf,
    style=\bf,
    headseparator={:},
    distance=0.5em,
]

% One extra line of space before and after every table
\setupfloat[table][
    spacebefore={medium,line},
    spaceafter={medium,line},
]

% =============================================================================
% CODE BLOCKS WITH SYNTAX HIGHLIGHTING (via vim)
% =============================================================================

% Load the vim module for syntax highlighting
\usemodule[vim]

% Python code
\definevimtyping[pythoncode][
    syntax=python,
    before={\blank[small]\startframedtext[
        frame=on,
        framecolor=codeborder,
        background=color,
        backgroundcolor=codebackground,
        width=broad,
        offset=0.5em,
    ]},
    after={\stopframedtext\blank[small]},
    style=\tt\small,
    numbering=no,
]

% Bash/shell code
\definevimtyping[bashcode][
    syntax=sh,
    before={\blank[small]\startframedtext[
        frame=on,
        framecolor=codeborder,
        background=color,
        backgroundcolor=codebackground,
        width=broad,
        offset=0.5em,
    ]},
    after={\stopframedtext\blank[small]},
    style=\tt\small,
    numbering=no,
]

% JSON code
\definevimtyping[jsoncode][
    syntax=json,
    before={\blank[small]\startframedtext[
        frame=on,
        framecolor=codeborder,
        background=color,
        backgroundcolor=codebackground,
        width=broad,
        offset=0.5em,
    ]},
    after={\stopframedtext\blank[small]},
    style=\tt\small,
    numbering=no,
]

% YAML code
\definevimtyping[yamlcode][
    syntax=yaml,
    before={\blank[small]\startframedtext[
        frame=on,
        framecolor=codeborder,
        background=color,
        backgroundcolor=codebackground,
        width=broad,
        offset=0.5em,
    ]},
    after={\stopframedtext\blank[small]},
    style=\tt\small,
    numbering=no,
]

% Plain text code (no highlighting)
\definetyping[textcode][
    before={\blank[small]\startframedtext[
        frame=on,
        framecolor=codeborder,
        background=color,
        backgroundcolor=codebackground,
        width=broad,
        offset=0.5em,
    ]},
    after={\stopframedtext\blank[small]},
    style=\tt\small,
]

% Inline-Code
\definetype[code][style=\tt]

% =============================================================================
% NOTICE BOXES
% =============================================================================

% Info box
\defineframedtext[infobox][
    frame=on,
    framecolor=primarycolor,
    background=color,
    backgroundcolor=infobackground,
    width=broad,
    offset=1em,
    before={\blank[medium]},
    after={\blank[medium]},
]

% Tip box (green)
\defineframedtext[tippbox][
    frame=on,
    framecolor=secondarycolor,
    background=color,
    backgroundcolor=secondarycolor:1,
    width=broad,
    offset=1em,
    before={\blank[medium]},
    after={\blank[medium]},
]

% Warning box (red)
\defineframedtext[warnbox][
    frame=on,
    framecolor=warningcolor,
    background=color,
    backgroundcolor=warningcolor:1,
    width=broad,
    offset=1em,
    before={\blank[medium]},
    after={\blank[medium]},
]
